\subsection*{2025}
\begin{enumerate}
	 \item Eder, J., Pühringer, S., \& Cserjan, L. (2025). Mobilitätswende produ­zieren: Warum die Bahn­industrie jetzt im Zentrum stehen muss.
	 \item \href{https://kompass.blog/public-brief/mobilitaetswende-produzieren-warum-die-bahnindustrie-jetzt-im-zentrum-stehen-muss/}{Porak, L., \& Hornykewycz, A. (2025). Mobilitätswende produzieren: Warum die Bahnindustrie jetzt im Zentrum stehen muss.}
\end{enumerate}
\subsection*{2024}
\begin{enumerate}
	 \item \href{https://www.awblog.at/Verteilung/Vermoegen-in-Oesterreich-stark-konzentriert-gering-besteuert}{Disslbacher, F., Hornykewycz, A., \& Kapeller, J. (2024). Vermögen in Österreich: stark konzen­triert, gering be­steuert.}
	 \item \href{https://www.awblog.at/Europa/Europa-braucht-industriepolitische-Antworten}{Porak, L. (2024). Globale geo-ökonomische Unordnung: Europa braucht industriepolitische Antworten.}
\end{enumerate}
\subsection*{2022}
\begin{enumerate}
	 \item \href{https://oxiblog.de/klima-markt-und-zukunftsbilder/}{Ötsch, W. (2022). Klima, Markt und Zukunftsbilder.}
\end{enumerate}
\subsection*{2021}
\begin{enumerate}
	 \item \href{https://www.awblog.at/Wirtschaft/budgetkuerzungen-durch-outputluecken-nonsens}{Heimberger, P. (2021). Budgetkürzungen durch „Outputlücken-Nonsens“.}
	 \item \href{https://www.awblog.at/Europa/europaeische-vermoegenssteuer-fuer-das-klima}{Kapeller, J., \& Wildauer, R. (2021). Eine europäische Vermögenssteuer für das Klima.}
	 \item \href{https://www.awblog.at/Wirtschaft/so-denken-oekonomen-ueber-wettbewerb}{Porak, L., Pühringer, S., \& Rath, J. (2021). So denken Ökonom*innen über Wettbewerb – eine Kritische Analyse des österreichischen Expert*innendiskurses.}
\end{enumerate}
\subsection*{2020}
\begin{enumerate}
	 \item \href{https://www.awblog.at/Allgemein/tuerkis-blaue-arbeitsmarkt-und-sozialpolitik-revisited}{Beyer, K., Griesser, M., \& Pühringer, S. (2020). Türkis-blaue Arbeitsmarkt- und Sozialpolitik revisited: zwischen Meritokratie und Wohlfahrtschauvinismus.}
	 \item \href{https://www.awblog.at/Wirtschaft/budgetpolitik-im-wirtschaftsabschwung-spielraeume}{Heimberger, P. (2020). Budgetpolitik im Wirtschaftsabschwung: erhebliche Spielräume vorhanden.}
	 \item \href{https://www.awblog.at/Europa/reform-der-budgetregeln-in-der-corona-krise-erforderlich}{Heimberger, P. (2020). Budgetpolitik in der Corona-Krise: Reform der Budgetregeln erforderlich.}
	 \item \href{https://www.awblog.at/Verteilung/oekonomische-globalisierung-und-einkommensungleichheit}{Heimberger, P. (2020). Wie die ökonomische Globalisierung die Einkommensungleichheit beeinflusst.}
\end{enumerate}
\subsection*{2019}
\begin{enumerate}
	 \item \href{https://www.awblog.at/Wirtschaft/einfluss-neoliberalismus-parteiprogramme}{Grimm, C. (2019). Der Einfluss des Neoliberalismus auf österreichische Parteiprogramme.}
	 \item \href{https://www.awblog.at/Arbeit/arbeitslosigkeit-in-europa}{Heimberger, P. (2019). Arbeitslosigkeit in Europa: Was man tun könnte.}
	 \item \href{https://www.awblog.at/Europa/italien-defizitverfahren-kontraproduktiv}{Heimberger, P. (2019). Italien vs. EU-Kommission: Warum ein Defizitverfahren kontraproduktiv wäre.}
	 \item \href{https://www.awblog.at/Europa/wirtschaftlichem-abschwung-entgegenwirken}{Heimberger, P., \& Pekanov, A. (2019). Dem wirtschaftlichen Abschwung entgegenwirken: Zur wichtigen Rolle der Fiskalpolitik.}
\end{enumerate}
\subsection*{2017}
\begin{enumerate}
	 \item \href{https://www.awblog.at/Bildung/soll-der-staat-bei-bildung-gesundheit-und-sozialem-kuerzen-aus}{Heimberger, P. (2017). Soll der Staat bei Bildung, Gesundheit und Sozialem kürzen? Austeritätspolitik seit der Finanzkrise im Vergleich.}
	 \item \href{https://www.awblog.at/Bildung/vorsicht-bei-laendervergleichen-insbesondere-bei-staatsausgabe}{Heimberger, P. (2017). Vorsicht bei Ländervergleichen – insbesondere bei Staatsausgaben!.}
	 \item \href{https://www.awblog.at/Allgemein/weniger-staatsausgaben-abbau-des-sozialstaats-und-vertiefung-v}{Heimberger, P. (2017). Weniger Staatsausgaben: Abbau des Sozialstaats und Vertiefung von Wirtschaftskrisen.}
	 \item \href{https://www.awblog.at/Europa/oesterreichs-bildungs-gesundheits-und-sozialausgaben-im-europa}{Heimberger, P. (2017). Österreichs Bildungs-, Gesundheits- und Sozialausgaben im europäischen Vergleich: Wenn der Staat spart, kann das für private Haushalte teuer werden.}
\end{enumerate}
\subsection*{2016}
\begin{enumerate}
	 \item \href{https://www.awblog.at/Allgemein/europaeische-schattenbanken-entwicklungen}{Beyer, K., \& Bräutigam, L. (2016). Das europäische Schattenbankensystem: Bestandsaufnahme und gegenwärtige Entwicklungen.}
	 \item \href{https://www.awblog.at/Wirtschaft/mehr-oeffentliche-investitionen-sind-sinnvoll-und-erforderlich}{Heimberger, P. (2016). Mehr öffentliche Investitionen sind sinnvoll und erforderlich.}
\end{enumerate}
\subsection*{2015}
\begin{enumerate}
	 \item \href{https://www.awblog.at/Allgemein/nachfrageseitige-ursachen-der-expansion-des-schattenbankensyst}{Beyer, K. (2015). Nachfrageseitige Ursachen der Expansion des Schattenbankensystems.}
	 \item \href{https://www.awblog.at/Allgemein/verteilungstene-perspektiven}{Kapeller, J., \& Schütz, B. (2015). Verteilungstendenzen im Kapitalismus: Globale Perspektiven.}
\end{enumerate}
\subsection*{2014}
\begin{enumerate}
	 \item \href{https://www.awblog.at/Europa/die-risiken-im-schatten-des-systems}{Beyer, K. (2014). Die Risiken im Schatten des Systems.}
	 \item \href{https://www.awblog.at/Europa/oekonomische-krisen-als-krankheiten-und-katastrophen}{Pühringer, S. (2014). Ökonomische Krisen als Krankheiten und Katastrophen?.}
\end{enumerate}
\subsection*{2013}
\begin{enumerate}
	 \item \href{https://www.awblog.at/Wirtschaft/ahnungslos-aber-nicht-tatenlos-wie-die-deutschsprachigen-okono}{Pühringer, S. (2013). Ahnungslos, aber nicht tatenlos – Wie ÖkonomInnen seit der Finanzkrise Politik mach(t)en.}
\end{enumerate}
